Section 3.3 shows that expanded training allows a model to read $r_t$
before predicting the later KCs of $q_t$, whereas all-in-one evaluation never
does. Therefore, the training process should limit the input infomation for each KC prediction such that the KC being predicted does not utilize information from other KCs within the same question. The resulting rule is called Pre-Question Conditioning (PQC). PQC constrains what the model may read; it is not a new architecture. Sections 4.2 and 4.3 show how to enforce it in a recurrent and in an attention backbone, and Section 4.4 states which parts come from prior work.

\subsection{The Rule}
\label{sec:rule-def}

Let $S_{<t}$ be the interactions of questions $q_1,\dots,q_{t-1}$, and let
$\mathcal{X}_{t,j}$ be what the model reads, besides $(q_t,c_{t,j})$, when it
predicts the $j$-th KC of $q_t$. Expanded training uses
\begin{equation}
  \mathcal{X}^{\mathrm{exp}}_{t,j}
  = S_{<t} \cup \{(q_t,c_{t,i},r_t)\}_{i<j},
  \label{eq:ctx-exp}
\end{equation}
while all-in-one evaluation uses
\begin{equation}
  \mathcal{X}^{\mathrm{eval}}_{t,j} = S_{<t}.
  \label{eq:ctx-eval}
\end{equation}
PQC requires, for every question $q_t$ and every $j \le m_t$:
\begin{enumerate}
  \item[(i)] During training, the label of a sentence is not included in the state used to predict the KC of that same question: $\mathcal{X}_{t,j}$ must not contain information
        $r_t$.
  \item[(ii)] All KCs of a question are predicted from the student’s state before that question, in line with the all-in-one protocol used for evaluation: all KCs of $q_t$ are predicted from
        the same context, the history before $q_t$, as in
        Eq.~\eqref{eq:ctx-eval}.
\end{enumerate}
Requirement (ii) implies requirement (i), so the two are nested. Requirement (i) alone prevents the model from simply copying the label. Requirement (ii) goes further and makes the training context exactly the same as the evaluation context. The two are stated separately because some implementations meet (i) but not (ii) (DKT-ML below), which makes it possible to measure the contribution of each step.
Under PQC we train with the per-KC binary cross-entropy
\begin{equation}
  \mathcal{L} = -\sum_{t}\sum_{j=1}^{m_t}
  \Big[ r_t \log \hat p_{t,j} + (1-r_t)\log(1-\hat p_{t,j}) \Big],
  \qquad
  \hat p_{t,j} = P\big(r_t = 1 \mid q_t, c_{t,j}, \mathcal{X}_{t,j}\big),
  \label{eq:loss}
\end{equation}
and evaluate exactly as in Section 3.2 (independent KC predictions, mean fusion).
Because $r_t$ is given to the model only after all $m_t$ predictions for $q_t$
have been made, the model cannot copy the preceding label. If every question has
a single KC, the rule changes nothing. It matters for training only when a
question has multiple KCs, either through its own rows or through the state it
leaves for later questions..

\subsection{Instantiation on Deep Knowledge Tracing}
Deep Knownledge Tracing (DKT) [cite] reads a sequence of tuple $(c, r)$, where $c$ is a KC and $r$ is the response to the question associated with KC $c$. Following Section 3.2, a question $q_t$ is expanded into a sequence of tuple ${(c_{t,1}, r_t), (c_{t,2}, r_t), ..., (c_{t, m_t}, r_t)}$. Let $h_{t, j}$ be LSTM state after reading $j$-th token of $q_t$, with $h_{t, 0} = h_{t-1, m_{t-1}}$ the state left by the previous question. Each tuple $(c_{t,j}, r_t)$ is encoded as a vector $\mathbf{x}_{t,j} \in \mathbb{R}^{d}$ and
\begin{equation}
  \mathbf{h}_{t,j} = \mathrm{LSTM}\big(\mathbf{x}_{t,j},\, \mathbf{h}_{t,j-1}\big), \qquad
  \mathbf{y}_{t,j} = \sigma\big(\mathbf{W}\mathbf{h}_{t,j} + \mathbf{b}\big).
\end{equation}

\paragraph{DKT (baseline)} Every token carries the true label and the next token is predicted from the previous one:
\begin{equation}
  \hat p_{t,j}=y_{t,j-1}[c_{t,j}],\qquad X_{t,j}=S_{<t}\cup\{(c_{t,i},r_t)\}_{i<j}.
\end{equation}
For $j>1$ the state has already read $r_t$, so neither (i) nor (ii) holds. This is the leakage of Section 3.3.

\paragraph{DKT-ML [cite].} The prediction rule is unchanged, but the label is hidden on every token except the last one of the question:
$$\rho_{t,j}=\begin{cases}MASK & j<m_t\\ r_t & j=m_t\end{cases}\qquad X_{t,j}=S_{<t}\cup\{(c_{t,i},\bot)\}_{i<j}.$$
Here $MASK$ is a mask symbol with its own input channel. Since $r_t$ is revealed only after the last KC, (i) holds. However, $c_{t,j}$ is still predicted after reading $c_{t,1},\dots,c_{t,j-1}$, so the context differs from evaluation and (ii) does not hold.

\paragraph{DKT-C (PQC, ours).} Tokens keep the true label, $\rho_{t,j}=r_t$, and are still read sequentially, but every KC of $q_t$ is predicted from the state before the question:
$$\hat p_{t,j}=y_{t-1,m_{t-1}}[c_{t,j}]=\sigma\!\big(Wh_{t,0}+b\big)[c_{t,j}],\qquad X_{t,j}=S_{<t}.$$
Neither $r_t$ nor the other KCs of $q_t$ is visible to the prediction of $c_{t,j}$, so (ii), and therefore (i), holds. The tokens of $q_t$ still pass through the LSTM, but they only update the state used for $q_{t+1}$.

\paragraph{DKT-ML-C.} This variant combines both changes: the labels are masked as in DKT-ML, $\rho_{t,j}=MASK$ for $j<m_t$, and every KC is predicted from $h_{t,0}$ as in DKT-C. The prediction context is the same as DKT-C, so (ii) holds. The two differ only in what the state stores for later questions: DKT-C keeps $r_t$ on every KC token, DKT-ML-C only on the last one.

\paragraph{DKT-Fuse.}
Instead of $m_t$ tokens, question $q_t$ is represented by a single token, obtained
by averaging the encodings of its KC tokens, and the LSTM advances by one step per
question:
\begin{equation}
  \bar{\mathbf{x}}_t = \frac{1}{m_t}\sum_{j=1}^{m_t}\mathbf{x}_{t,j},
  \qquad
  \mathbf{h}_t = \mathrm{LSTM}\big(\bar{\mathbf{x}}_t,\, \mathbf{h}_{t-1}\big),
  \qquad
  \mathbf{y}_t = \sigma\big(\mathbf{W}\mathbf{h}_t + \mathbf{b}\big).
\end{equation}
Every KC of $q_t$ is predicted from the state before the question,
\begin{equation}
  \hat p_{t,j} = \mathbf{y}_{t-1}[c_{t,j}], \qquad X_{t,j} = S_{<t},
\end{equation}
so neither $r_t$ nor the other KCs of $q_t$ is visible, and (i) and (ii) hold. The
original fusion approach of Badran and Preisach~\cite{badran} averages the KC
predictions into a question-level probability,
$\hat p_t=\frac{1}{m_t}\sum_{j=1}^{m_t}\hat p_{t,j}$, and applies the binary
cross-entropy to $\hat p_t$. We do not average before the loss: each KC receives its
own term of Eq.~(4). Averaging is applied only at evaluation, as in
Section~3.2.

\subsection{Design Space and Provenance}
\label{sec:rule-origin}

Table~\ref{tab:design} lists all variants. The principle of hiding same-question
interactions in attention was introduced as AKT-QM by Badran and
Preisach~\cite{badran2025}; masked labels (DKT-ML) and fused updates
(DKT-Fuse) also come from their work~\cite{badran2025,badran2026}. We do not
claim these mechanisms. What we add is (a) stating the rule as two separate
requirements, (i) and (ii), (b) DKT-PQC-Seq, which satisfies both without extra
parameters, and DKT-PQC-Mask as an intermediate step, and (c) a controlled
measurement of each step, and of the rule on simpleKT, under the all-in-one
protocol of pyKT. Question-level models such as QIKT build one step per
question, so they have no intra-question leakage by construction and are not
the target of the rule. A further control, which changes only checkpoint
selection, is described in Sect.~\ref{sec:setup}.

\begin{table}[t]
  \caption{Design space of the training rule. $h_{t,j-1}$ is the state after the
  first $j-1$ KCs of $q_t$ and $h_{t-1}$ the state before $q_t$; (i) and (ii) are
  the requirements of Sect.~\ref{sec:rule-def}. Parameter counts are for our
  XES3G5M configuration.}
  \label{tab:design}
  \centering
  \scriptsize
  \setlength{\tabcolsep}{3pt}
  \begin{tabular}{@{}l p{2.6cm} p{1.5cm} c c r p{1.9cm}@{}}
    \hline
    Variant & Update by the KCs of $q_t$ & KC $j$ read from & (i) & (ii) & Params & Origin \\
    \hline
    DKT-Exp          & one step per KC, true label & $h_{t,j-1}$ & no  & no  & 588,465   & \cite{piech2015dkt,pykt} \\
    DKT-ML           & one step per KC, MASK except last & $h_{t,j-1}$ & yes & no  & 674,965   & \cite{badran2025} \\
    DKT-PQC-Mask     & as DKT-ML & $h_{t-1}$ & yes & yes & 674,965   & ours \\
    DKT-PQC-Seq      & one step per KC, true label & $h_{t-1}$ & yes & yes & 588,465   & ours \\
    DKT-Fuse         & one step per question & $h_{t-1}$ & yes & yes & 588,465   & \cite{badran2025}; per-KC loss ours \\
    simpleKT         & attention over rows $j<i$ & rows $j<i$ & no  & no  & 3,754,497 & \cite{liu2023simplekt,pykt} \\
    simpleKT-QM      & attention over rows $g_j<g_i$ & rows $g_j<g_i$ & yes & yes & 3,754,497 & AKT-QM principle~\cite{badran2025} \\
    \hline
  \end{tabular}
  % NOTE: regenerate the Params column by script from registry.csv; simpleKT-QM count to be confirmed (same as simpleKT).
\end{table}